\documentclass[11pt]{article}
\usepackage[utf8]{inputenc}
\usepackage[margin=1in]{geometry}
\usepackage{amsmath}
\usepackage{graphicx}
\usepackage{booktabs}
\usepackage{hyperref}
\usepackage{natbib}
\usepackage{setspace}
\onehalfspacing

\title{Antipsychotic Drugs Selectively Decorrelate Long-Range Interactions in Deep Cortical Layers}
\author{Author A$^{1}$, Author B$^{1,2}$, Author C$^{2}$ \\
\small $^{1}$Department of Neuroscience, University of Research \\
\small $^{2}$Institute for Brain and Behavior}
\date{}

\begin{document}
\maketitle

\begin{abstract}
Antipsychotic drugs are the primary pharmacological treatment for psychosis, yet their circuit-level effects in the cortex remain poorly understood. Psychosis is thought to involve impaired distinction between externally driven and self-generated neural activity, a computation that may rely on predictive processing across cortical layers. Here we used widefield calcium imaging through crystal skull cranial windows in head-fixed mice to map drug effects across genetically defined cortical neuron types. We found that layer 5 intratelencephalic (L5 IT, Tlx3-positive) neurons most strongly distinguished closed-loop (locomotion drives visual flow) from open-loop (replayed visual flow) conditions compared to superficial layer neurons. The atypical antipsychotic clozapine selectively reduced inter-regional correlations in L5 IT neurons without affecting superficial layers. This deep-layer decorrelation was replicated by two additional antipsychotics---aripiprazole and haloperidol---but not by the psychostimulant amphetamine. Antipsychotic treatment preferentially reduced negative prediction error responses and their propagation from V1 to higher visual areas in L5 IT neurons. These results reveal a layer-specific cortical mechanism of antipsychotic drug action centered on L5 IT neurons that mediate long-range cortical communication and visuomotor integration.
\end{abstract}

\section{Introduction}

Psychotic disorders, including schizophrenia, involve a fundamental disruption in the brain's ability to distinguish between externally generated sensory input and internally generated predictions \citep{friston2005,adams2013}. This impairment may arise from dysfunctional predictive processing, in which the brain generates expectations about incoming sensory information and computes prediction errors when expectations are violated \citep{rao1999,keller2018}. The predictive processing framework proposes a laminar organization in which superficial cortical layers (L2/3) signal prediction errors while deep layers (L5/6) carry internal model predictions and mediate long-range feedback communication \citep{bastos2012}.

Antipsychotic drugs are the primary treatment for psychosis and act broadly on multiple neuromodulatory receptors, including dopamine D2, serotonin 5-HT2A, muscarinic, and histaminergic receptors \citep{kapur2001,meltzer2013}. Despite decades of clinical use, the cortical circuit-level mechanisms of antipsychotic action remain unclear. Prior neuroimaging studies in humans have shown widespread effects on functional connectivity, but these studies lack the cell-type and layer specificity needed to identify the precise circuit elements affected \citep{honey2004}. Whether antipsychotics act preferentially on specific cortical layers or cell types has not been systematically investigated.

Recent advances in transgenic mouse lines expressing genetically encoded calcium indicators in specific cell populations enable the study of drug effects with cell-type resolution across the entire dorsal cortex \citep{dana2019,madisen2015}. Widefield calcium imaging through crystal skull or clear skull preparations provides large-scale access to cortical activity with sufficient spatial resolution to distinguish between cortical regions \citep{mohajerani2013,musall2019}. By combining these tools with virtual reality paradigms that manipulate the relationship between motor output and sensory feedback, it is possible to probe visuomotor integration and prediction error computation at the cell-type level \citep{keller2012}.

Here we used widefield calcium imaging through crystal skull cranial windows in mice expressing GCaMP6 in specific cell types via Cre-driver lines to systematically map the effects of antipsychotic drugs across cortical layers and neuron types. We compared closed-loop locomotion (where running on a spherical treadmill drives visual flow through a virtual corridor) with open-loop conditions (where previously recorded visual flow is replayed) to assess visuomotor integration. We found that L5 IT neurons (labeled by Tlx3-Cre) most strongly distinguished these conditions and were selectively affected by antipsychotic treatment.

\section{Methods}

\subsection{Animals and Transgenic Lines}

All procedures were approved by the institutional animal care and use committee. Nine transgenic mouse lines were used (Table~\ref{tab:mice}), including a brain-wide GCaMP expression line (C57BL/6 with AAV-PHP.eB) and eight Cre-driver lines crossed with the Ai148 GCaMP6 reporter to achieve cell type-specific expression.

\begin{table}[h]
\centering
\caption{Mouse lines and sample sizes used in this study.}
\label{tab:mice}
\begin{tabular}{llc}
\toprule
Genotype & Target Cell Type & $n$ (mice) \\
\midrule
C57BL/6 + AAV-PHP.eB & Brain-wide & 6 \\
Emx1-Cre & Cortical excitatory & 4 \\
Cux2-CreERT2 & L2/3 and L4 excitatory & 4 \\
Scnn1a-Cre & L4 excitatory subset & 7 \\
Tlx3-Cre $\times$ Ai148 & L5 IT neurons & 25 \\
Ntsr1-Cre & L6 excitatory subset & 3 \\
PV-Cre & PV interneurons & 2 \\
VIP-Cre & VIP interneurons & 6 \\
SST-Cre & SST interneurons & 5 \\
\bottomrule
\end{tabular}
\end{table}

\subsection{Surgical Preparation and Imaging}

Mice received crystal skull cranial windows for chronic widefield imaging \citep{kim2016}. Under isoflurane anesthesia, the skull was replaced with a transparent crystal skull prosthesis providing optical access to the entire dorsal cortex. After recovery, mice were head-fixed on a spherical treadmill positioned inside a virtual corridor with vertical black and white bars. GCaMP6 was excited with 470~nm LED illumination and an isosbestic wavelength (405~nm) was used for hemodynamic correction. Fluorescence was captured with an sCMOS camera. The signal attenuation half-depth was approximately 100~$\mu$m, favoring detection of superficial activity in brain-wide preparations and layer-specific activity in Cre-line preparations.

\subsection{Behavioral Paradigm}

Mice navigated a virtual corridor in two conditions: (1) closed-loop, where locomotion speed directly controlled visual flow speed, and (2) open-loop, where visual flow from a previous closed-loop session was replayed independent of the animal's locomotion. Locomotion onsets were identified as transitions from stationary to running and were aligned across conditions for comparison.

\subsection{Drug Administration}

Single intraperitoneal injections of the following drugs were administered: clozapine (atypical antipsychotic), aripiprazole (atypical antipsychotic, D2 partial agonist), haloperidol (typical antipsychotic, D2 antagonist), and amphetamine (psychostimulant, DA/NE releaser). Each animal served as its own control with pre-drug and post-drug imaging sessions performed on the same day.

\subsection{Data Analysis}

Twelve bilateral dorsal cortical regions of interest (ROIs) were defined. Inter-regional correlations were computed for each pair of ROIs during locomotion bouts. Distance-dependent correlation analyses related inter-regional correlation magnitude to the physical distance between ROIs, normalized by bregma-lambda distance. Statistical comparisons used hierarchical bootstrap with 90\% confidence intervals to account for the nested structure of the data (trials within sessions within mice).

\section{Results}

\subsection{L5 IT Neurons Best Distinguish Visuomotor Conditions}

We first compared cortical activation patterns during closed-loop and open-loop locomotion onsets across the eight Cre lines (Table~\ref{tab:closedopen}). Superficial layer neurons (Cux2, L2/3-L4; Scnn1a, L4) showed relatively weak distinctions between the two conditions. In contrast, Tlx3-positive L5 IT neurons showed the strongest distinction, with markedly different activation patterns depending on whether locomotion was coupled to visual flow. Ntsr1-positive L6 neurons showed a similar level of distinction. Among inhibitory neuron types, PV, VIP, and SST interneurons showed moderate distinctions.

\begin{table}[h]
\centering
\caption{Closed-loop versus open-loop locomotion onset distinction by cell type.}
\label{tab:closedopen}
\begin{tabular}{llc}
\toprule
Neuron Type & Layer & Distinction Strength \\
\midrule
Cux2 & L2/3-L4 & Weak \\
Scnn1a & L4 & Weak \\
Tlx3 & L5 IT & Strongest \\
Ntsr1 & L6 & Similar to Tlx3 \\
PV & Mixed & Moderate \\
VIP & Mixed & Moderate \\
SST & Mixed & Moderate \\
\bottomrule
\end{tabular}
\end{table}

These findings indicate that deep cortical layers, particularly L5 IT neurons, carry the most information about whether visual input is self-generated (closed-loop) or externally imposed (open-loop). This is consistent with the predictive processing framework, which assigns internal model representations to deep cortical layers.

\subsection{Visuomotor Prediction Error Responses}

We next examined responses to visuomotor mismatches (negative prediction errors) by briefly halting visual flow during closed-loop locomotion. In brain-wide imaging, mismatch stimuli evoked strong activation originating in V1 that propagated to higher visual areas. Drifting gratings, which provide visual input without motor prediction, produced a distinct activation pattern. In Tlx3 L5 IT recordings, mismatch responses were robust in V1 and propagated strongly to V2am, consistent with a role for L5 IT neurons in transmitting prediction error signals between cortical areas.

\subsection{Clozapine Selectively Decorrelates L5 IT Inter-Regional Activity}

Administration of clozapine produced a dramatic reduction in inter-regional correlations that was highly specific to L5 IT neurons (Table~\ref{tab:clozapine}). In brain-wide (C57BL/6) imaging, clozapine caused a moderate reduction in the correlation matrix. However, when the same analysis was performed in Tlx3 L5 IT neurons, the decorrelation was far more pronounced. Critically, correlations in superficial layer neurons (Cux2, Scnn1a) were minimally affected by clozapine.

\begin{table}[h]
\centering
\caption{Clozapine effects on inter-regional correlations by cell type.}
\label{tab:clozapine}
\begin{tabular}{lccc}
\toprule
Cell Type & $n$ & Post-Clozapine Change & Effect Magnitude \\
\midrule
Brain-wide (C57BL/6) & 4 & Reduced & Moderate \\
Tlx3 (L5 IT) & 5 & Strongly reduced & Dramatic \\
Cux2 (L2/3-L4) & 4 & Minimal change & Weak \\
Scnn1a (L4) & 7 & Minimal change & Weak \\
\bottomrule
\end{tabular}
\end{table}

The decorrelation was distance-dependent: long-range correlations (between distant cortical areas) were more strongly reduced than short-range correlations. This pattern suggests that clozapine specifically disrupts the long-range communication mediated by L5 IT neurons, which are known to send long-range intracortical projections.

\subsection{Multiple Antipsychotics Share L5 IT Decorrelation; Amphetamine Does Not}

To determine whether the L5 IT decorrelation is a general property of antipsychotic drugs, we tested aripiprazole (atypical, D2 partial agonist) and haloperidol (typical, D2 antagonist). Both drugs reproduced the strong reduction in L5 IT inter-regional correlations observed with clozapine. Despite their different receptor binding profiles, all three antipsychotics converged on the same circuit-level effect.

In contrast, amphetamine, a psychostimulant that increases dopamine and norepinephrine release, did not produce significant decorrelation in L5 IT neurons. This dissociation between antipsychotics and amphetamine is particularly notable because amphetamine can induce psychotic symptoms in humans, suggesting that the pro-psychotic and anti-psychotic effects of these drugs may operate through distinct circuit mechanisms.

\subsection{Antipsychotics Reduce Prediction Error Responses in L5 IT}

Antipsychotic treatment preferentially reduced negative prediction error (mismatch) responses in V1 L5 IT neurons. The amplitude of mismatch-evoked activity was reduced, and critically, the propagation of mismatch signals from V1 to V2am was impaired. In contrast, responses to drifting gratings were less affected by antipsychotic treatment, suggesting a preferential impact on prediction error computation rather than general visual responsiveness.

This finding connects the inter-regional decorrelation effect to a specific computational function: antipsychotics disrupt the transmission of prediction error signals through L5 IT neurons, which may normalize the aberrant prediction error processing thought to underlie psychotic symptoms.

\subsection{Hemodynamic Controls}

To ensure that the observed effects reflected genuine neural activity changes rather than hemodynamic artifacts, we performed several control experiments. Crystal skull preparations showed smaller hemodynamic artifacts than clear skull preparations. Imaging in eGFP-expressing mice (no GCaMP) confirmed that hemodynamic signals alone could not account for the observed drug effects. Isosbestic (405~nm) correction was applied but provided only incomplete hemodynamic correction, consistent with the known limitations of single-wavelength correction methods.

\section{Discussion}

Our findings reveal a layer-specific cortical mechanism of antipsychotic drug action. Three antipsychotic drugs with different receptor binding profiles---clozapine, aripiprazole, and haloperidol---all selectively reduced long-range inter-regional correlations in L5 IT neurons without substantially affecting superficial cortical layers. This effect was not shared by the psychostimulant amphetamine, suggesting it is specific to the therapeutic mechanism of antipsychotic action rather than a general consequence of neuromodulatory manipulation.

The preferential targeting of L5 IT neurons by antipsychotics is consistent with the known biology of these neurons. L5 IT neurons project long-range within the cortex and to subcortical targets, making them key mediators of inter-regional communication \citep{harris2019a,baker2018}. They express high levels of dopamine and serotonin receptors, the primary targets of antipsychotic drugs \citep{gaspar2009}. Their position in the cortical circuit---receiving inputs from all layers and projecting widely---makes them ideal candidates for integrating prediction signals and transmitting them across areas.

The connection between L5 IT decorrelation and reduced prediction error propagation provides a circuit-level explanation for antipsychotic efficacy. If psychosis involves aberrant prediction error signals that are transmitted through L5 IT neurons, then reducing the gain of this transmission channel would normalize sensory processing and reduce hallucinatory and delusional experiences \citep{corlett2019}. Our data show that antipsychotics specifically reduce the propagation of mismatch (prediction error) signals from V1 to V2am in L5 IT neurons, providing direct evidence for this mechanism.

The lack of effect in superficial layers is noteworthy. In the predictive processing framework, superficial layers are thought to carry prediction error signals upward through the cortical hierarchy \citep{bastos2012}. Our finding that antipsychotics primarily affect deep layers suggests that the therapeutic mechanism may involve modulating the internal model or top-down feedback, rather than directly dampening prediction errors at their source.

Several limitations should be considered. First, widefield calcium imaging measures population-level activity and cannot resolve individual neurons. Second, the drugs were administered systemically, so we cannot distinguish direct cortical effects from indirect effects mediated through subcortical circuits. Third, the mouse virtual reality paradigm, while providing controlled visuomotor conditions, may not fully capture the complexity of sensory processing relevant to human psychosis \citep{keller2012}. Future studies using two-photon imaging with cell-type-specific indicators could provide single-neuron resolution of antipsychotic effects in deep cortical layers.

\section{Conclusion}

We have demonstrated that antipsychotic drugs selectively decorrelate long-range inter-regional activity in L5 IT cortical neurons, a cell type critical for visuomotor integration and prediction error propagation. This layer-specific mechanism is shared across typical and atypical antipsychotics but not by the psychostimulant amphetamine, suggesting it represents a core feature of antipsychotic therapeutic action. These findings provide a cell type-resolved framework for understanding how antipsychotic drugs normalize cortical processing and point to L5 IT neurons as a key circuit element in the pathophysiology and treatment of psychosis.

\bibliographystyle{plainnat}
\bibliography{references}

\end{document}
